\section{1. O que são esses materiais?}

\begin{frame}{Nem conduz, nem isola}
  \begin{columns}[T]
    \begin{column}{0.52\linewidth}
      \begin{itemize}
        \item No \termo{cobre} os elétrons estão soltos o tempo todo, então
              ele conduz sempre.
        \item No \termo{vidro} todos os elétrons estão presos nas ligações,
              então ele nunca conduz.
        \item No \termo{silício} os elétrons estão presos, mas soltam com
              pouca energia. Calor, luz ou uma tensão já bastam.
      \end{itemize}
      \vspace{0.2cm}
      \centering
      É o único dos três em que \alert{nós escolhemos o quanto ele conduz}.
    \end{column}
    \begin{column}{0.44\linewidth}
      \centering
      \begin{tikzpicture}[scale=0.68,transform shape]
        % cobre: as duas faixas se encostam, não há degrau
        \fill[mPrimary!30] (0,0) rectangle (1.6,1.1);
        \node[font=\scriptsize] at (0.8,0.55) {presos};
        \fill[mAccent!30] (0,1.1) rectangle (1.6,2.2);
        \node[font=\scriptsize] at (0.8,1.65) {soltos};
        \node[font=\scriptsize\bfseries] at (0.8,-0.45) {Cobre};
        \node[font=\tiny,text=mNeutral] at (0.8,-0.95) {sem degrau};

        \foreach \x/\gap/\nome/\tam in {3/1.1/Silício/pequeno, 6/2.6/Vidro/grande} {
          \fill[mPrimary!30] (\x,0) rectangle (\x+1.6,1.1);
          \node[font=\scriptsize] at (\x+0.8,0.55) {presos};
          \fill[mAccent!30] (\x,1.1+\gap) rectangle (\x+1.6,2.2+\gap);
          \node[font=\scriptsize] at (\x+0.8,1.65+\gap) {soltos};
          \draw[<->,thick,mNeutral] (\x+0.8,1.1) -- (\x+0.8,1.1+\gap);
          \node[font=\scriptsize\bfseries] at (\x+0.8,-0.45) {\nome};
          \node[font=\tiny,text=mNeutral] at (\x+0.8,-0.95) {degrau \tam};
        }
      \end{tikzpicture}
      \\[0.2cm]
      {\scriptsize O tamanho do degrau que o elétron precisa vencer é o
      que separa os três casos.}
    \end{column}
  \end{columns}
  \fonte{Callister, \emph{Ciência e Engenharia de Materiais}, cap.~18.}
\end{frame}

\begin{frame}{Sozinho o silício não serve}
  \begin{itemize}
    \item Em silício puro, só 1 átomo a cada 5 trilhões solta um elétron na
          temperatura ambiente. Conduz tão pouco que não dá para fazer nada
          com ele.
    \item A solução é sujar o cristal \alert{de propósito}, trocando um átomo
          a cada milhão por um vizinho da tabela periódica. Isso se chama
          \termo{dopagem}.
  \end{itemize}
  \vspace{0.2cm}
  \begin{columns}[T]
    \begin{column}{0.48\linewidth}
      \begin{block}{Fósforo, tipo N}
        \small Tem 5 elétrons para ligar e o silício só precisa de 4. O que
        sobra fica \termo{livre} para conduzir.
      \end{block}
    \end{column}
    \begin{column}{0.48\linewidth}
      \begin{block}{Boro, tipo P}
        \small Tem só 3 elétrons. A ligação que fica faltando é uma
        \termo{lacuna}, e ela também caminha pelo cristal.
      \end{block}
    \end{column}
  \end{columns}
  \vspace{0.2cm}
  \centering\small
  Uma impureza em um milhão faz a condutividade subir \alert{cem mil vezes}.
\end{frame}

\begin{frame}{Quem é quem na família}
  \begin{table}
    \centering
    \begin{tabular}{llc l}
      \toprule
      \textbf{Material} & & \textbf{Degrau (eV)} & \textbf{Onde aparece} \\
      \midrule
      Silício            & Si   & 1,1 & Chips, memórias, quase tudo \\
      Arseneto de gálio  & GaAs & 1,4 & Amplificadores de celular, LED \\
      Nitreto de gálio   & GaN  & 3,4 & Antenas 5G, carregadores \\
      Fosfeto de índio   & InP  & 1,3 & Lasers de fibra óptica \\
      \bottomrule
    \end{tabular}
  \end{table}
  \begin{itemize}\small
    \item Degrau pequeno solta elétron fácil, mas vaza corrente e esquenta.
    \item Degrau grande aguenta tensão alta, e por isso o GaN vai para as
          antenas 5G.
  \end{itemize}
  \gancho{As duas bases são \termo{silício e gálio}. O silício domina por ser
  barato; o gálio entra onde é preciso alta frequência ou luz.}
\end{frame}
